\pagestyle{empty}
\tight
\begin{tblr}{width=\textwidth,colspec={|Q[c,m,1.9cm]|X[l,m]|},hlines,vlines,hspan=minimal,row{1}={bg=headblue},rowsep=1pt,colsep=3pt}
\SetCell{c,m}\hfil\logo\hfil & \textbf{POLITEKNIK NEGERI MALANG}\par\textbf{JURUSAN TEKNOLOGI INFORMASI}\par\textbf{PROGRAM STUDI: D4 TEKNIK INFORMATIKA} \\
\SetCell[c=2]{c} \textbf{RENCANA PEMBELAJARAN SEMESTER (RPS)} & \\
\end{tblr}
\begin{longtblr}{width=\textwidth,colspec={|Q[l,m,1.9cm]|X[0.85,l,m]|X[1.45,l,m]|X[0.75,l,m]|X[1.15,l,m]|},hlines,vlines,hspan=minimal,rowsep=1pt,colsep=2pt,row{1,3}={bg=headgray,font=\bfseries}}
MATA KULIAH & KODE & BOBOT (SKS)/jam & SEMESTER & TGL. PENYUSUNAN \\
Konsep Teknologi Informasi & RTI251002 & 2 SKS/4 jam & 1 & 20 Agustus 2025 \\
OTORISASI & Dosen Pengembang RPS & Koordinator RMK & \SetCell[c=2]{l} Ka PRODI &  \\
 & Ariadi Retno Ririd, S.Kom., M.Kom. & Ariadi Retno Ririd, S.Kom., M.Kom. & \SetCell[c=2]{l}{\vspace{8mm}\par Dr. Ely Setyo Astuti, S.T., M.T.} &  \\
\SetCell[r=14]{l} Capaian Pembelajaran (CP) & \SetCell[c=4]{l,bg=headgray,font=\bfseries} Capaian Pembelajaran Lulusan Program Studi (CPL-Prodi) &  &  &  \\
 & CPL02 & \SetCell[c=3]{l} Mampu menerapkan prinsip rekayasa sistem dan perangkat lunak untuk menghasilkan solusi aplikasi multi-platform sesuai kebutuhan. &  &  \\
 & CPL05 & \SetCell[c=3]{l} Mampu menghasilkan solusi teknologi komputasi dan infrastruktur digital yang sesuai dengan kebutuhan organisasi dan pengguna. &  &  \\
 & CPL09 & \SetCell[c=3]{l} Mampu menjelaskan cara kerja sistem komputer dan menerapkan pendekatan komputasi untuk menyelesaikan masalah. &  &  \\
 & \SetCell[c=4]{l,bg=headgray,font=\bfseries} Capaian Pembelajaran mata kuliah (CPL-MK) &  &  &  \\
 & CPMK0208 & \SetCell[c=3]{l} Mahasiswa mampu memilih arsitektur computing system sesuai kebutuhan aplikasi multi-platform. &  &  \\
 & CPMK0901 & \SetCell[c=3]{l} Mahasiswa mampu menjelaskan cara kerja sistem komputer dan komponen utamanya dalam penyelesaian masalah komputasi. &  &  \\
 & CPMK0502 & \SetCell[c=3]{l} Mahasiswa mampu mengidentifikasi dan menangani isu keamanan pada sistem operasi, jaringan komputer, dan platform cloud. &  &  \\
 & \SetCell[c=4]{l,bg=headgray,font=\bfseries} Kemampuan akhir tiap tahapan belajar (Sub-CPMK) &  &  &  \\
 & SCPMK0901-00201 & \SetCell[c=3]{l} Mahasiswa mampu menjelaskan konsep teknologi informasi, inovasi teknologi, perkembangan IPTEK dan ICT, serta etika rekayasa di bidang TI. &  &  \\
 & SCPMK0901-00202 & \SetCell[c=3]{l} Mahasiswa mampu menjelaskan cara kerja sistem komputer beserta struktur, arsitektur, dan komponen utamanya. &  &  \\
 & SCPMK0901-00203 & \SetCell[c=3]{l} Mahasiswa mampu menerapkan representasi data, aljabar Boolean, dan flowchart untuk menyelesaikan masalah komputasi sederhana. &  &  \\
 & SCPMK0208-00204 & \SetCell[c=3]{l} Mahasiswa mampu memilih jaringan komputer, aplikasi TI, dan sertifikasi bidang TI yang sesuai dengan kebutuhan pengguna dan organisasi. &  &  \\
 & SCPMK0502-00205 & \SetCell[c=3]{l} Mahasiswa mampu mengidentifikasi aspek keamanan dasar pada sistem operasi, jaringan komputer, dan platform cloud sebagai bagian dari wawasan teknologi informasi. &  &  \\
\SetCell[r=2]{l} Penilaian & \SetCell[c=4]{l} *Nilai CPMK didapat dari agregasi nilai SUB-CPMK yang dibebankan &  &  &  \\
 & \SetCell[c=4]{l}{\tiny\begin{tblr}{width=\linewidth,colspec={|Q[c,m,0.1\linewidth]|Q[c,m,0.16\linewidth]|X[l,m]|X[l,m]|X[l,m]|X[l,m]|X[l,m]|Q[c,m,0.08\linewidth]|},hlines,vlines,hspan=minimal,rowsep=1pt,colsep=0.7pt,row{1,2}={bg=headgray,font=\bfseries}}
Kode CPMK & Kode SCPMK & \SetCell[c=5]{c} Bobot per Bentuk Penilaian & & & & & Total Bobot \\
 & & Tugas & Kuis 1 & UTS & Kuis 2 & UAS & \\
\SetCell[r=1]{c} \SetCell[r=3]{c} CPMK0901 & SCPMK0901-00201 & 10\%: rangkuman konsep TI, inovasi, IPTEK, etika, ICT & 12\%: konsep dasar TI & 7\%: konsep TI, IPTEK, dan etika & & & 29\% \\
\SetCell[r=2]{c}  & SCPMK0901-00202 & 4\%: ulasan sistem komputer & & 8\%: struktur dan arsitektur sistem komputer & & & 12\% \\
 & SCPMK0901-00203 & 6\%: latihan representasi data, Boolean, flowchart & & & 12\%: komputasi dan logika & 15\%: penerapan komputasi & 33\% \\
\SetCell[r=1]{c} \SetCell[r=1]{c} CPMK0208 & SCPMK0208-00204 & 6\%: ulasan jaringan, aplikasi TI, sertifikasi & & & & 17\%: jaringan dan aplikasi TI & 23\% \\
CPMK0502 & SCPMK0502-00205 & & & & & 3\%: aspek keamanan dasar & 3\% \\
\SetCell[c=7]{r}\textbf{Total Per Penilaian} & & & & & & & \textbf{100\%} \\
\end{tblr}} &  &  &  \\
Diskripsi Singkat Mata Kuliah & \SetCell[c=4]{l} Dalam mata kuliah ini akan dibahas tentang konsep teknologi, inovasi teknologi, perkembangan IPTEK, etika rekayasa, perkembangan ICT, sistem komputer, konsep sistem komputer, representasi data, aljabar Boolean, flowchart, jaringan komputer dan internet, aplikasi TI di berbagai bidang, serta sertifikasi bidang TI. &  &  &  \\
Bahan Kajian / Materi Pembelajaran & \SetCell[c=4]{l}\begin{enumerate}\item Konsep teknologi dan inovasi teknologi\item Perkembangan IPTEK dan etika rekayasa\item Perkembangan ICT\item Sistem komputer dan konsep sistem komputer\item Representasi data, aljabar Boolean, dan flowchart\item Jaringan komputer dan internet\item Aplikasi TI di berbagai bidang dan sertifikasi bidang TI\end{enumerate} &  &  &  \\
\SetCell[r=2]{l} Pustaka & \SetCell[c=4]{l} \textbf{Utama:}\begin{enumerate}\item William Stallings. 2007. Data and Computer Communications. 8th Edition. Pearson Education, Inc. Pearson Prentice Hall.\end{enumerate} &  &  &  \\
 & \SetCell[c=4]{l} \textbf{Pendukung:}\begin{enumerate}\item Davis, W.S. Computers and Information Systems: An Introduction. West Publishing Company.\item Situs web dan artikel KTI.\end{enumerate} &  &  &  \\
Media Pembelajaran & \SetCell[c=2]{l} \textbf{Software:} LMS, browser, office suite, aplikasi ujian online. &  & \SetCell[c=2]{l} \textbf{Hardware:} Komputer/laptop, proyektor. &  \\
Nama Dosen Pengampu & \SetCell[c=4]{l}\begin{enumerate}\item Ariadi Retno Ririd, S.Kom., M.Kom.\item Budi Harijanto, S.T., M.MKom.\item Deddy Kusbianto PA, Ir., M.MKom.\item Banni Satria Andoko, S.Kom., M.MSI., Dr.Eng.\end{enumerate} &  &  &  \\
Matakuliah Syarat & \SetCell[c=4]{l} - &  &  &  \\
\end{longtblr}
\begin{longtblr}{width=\textwidth,colspec={|Q[c,m,0.06\textwidth]|X[1.35,l,m]|X[1.08,l,m]|X[1.16,l,m]|X[0.7,l,m]|X[1.24,l,m]|X[1.06,l,m]|X[0.98,l,m]|Q[c,m,0.05\textwidth]|},hlines,vlines,hspan=minimal,rowhead=2,row{1,2}={bg=headgreen,font=\bfseries},rowsep=1pt,colsep=1.5pt}
Mgg.\par Ke & Kemampuan Akhir Yang Direncanakan (Sub-CP-MK) & Bahan kajian (Materi Pembelajaran) & Bentuk dan Metode Pembelajaran & Estimasi Waktu & Pengalaman Belajar Mahasiswa & Kriteria \& Bentuk Penilaian & Indikator Penilaian & Bobot Penilaian (\%) \\
(1) & (2) & (3) & (4) & (5) & (6) & (7) & (8) & (9) \\
1 & SCPMK0901-00201 & Konsep teknologi informasi: fungsi, pemanfaatan, pengelompokan, komponen pembangun, peranan dasar, dan trend TI. & Bentuk: Daring (Online). Metode: Contextual Teaching and Learning (CTL) dan Self Directed Learning (SDL). Penugasan: tugas individu merangkum ulasan teknologi informasi. & 1 x 3 x 50' tatap muka. & Mahasiswa mempelajari konsep teknologi informasi serta memahami peranan dan fungsi TI dalam dunia saat ini. & Kriteria: ketepatan dan penguasaan. Bentuk: keaktifan diskusi kelompok serta kesesuaian, kejelasan, dan kelugasan menjawab tugas. & Mahasiswa dapat menjelaskan konsep TI, fungsi dan pemanfaatan TI, pengelompokan TI, komponen pembangun TI, dan peranan dasar TI. & 2 \\
2 & SCPMK0901-00201 & Inovasi teknologi: perbedaan inovasi sistem informasi dan teknologi informasi modern beserta contohnya. & Bentuk: Daring (Online). Metode: Self Directed Learning (SDL) dan Cooperative Learning (CoL). Penugasan: tugas kelompok merangkum inovasi TI era modern. & 1 x 3 x 50' tatap muka. & Mahasiswa mempelajari inovasi teknologi informasi saat ini dan contoh-contoh inovasi yang berkembang. & Kriteria: ketepatan dan penguasaan. Bentuk: keaktifan diskusi kelompok serta kesesuaian, kejelasan, dan kelugasan menjawab tugas. & Mahasiswa mampu memahami dan menjelaskan perbedaan inovasi sistem informasi dan teknologi informasi modern. & 2 \\
3 & SCPMK0901-00201 & Pengertian dan perkembangan IPTEK, perkembangan IPTEK dalam bidang pendidikan, serta pengaruh IPTEK dan solusinya. & Bentuk: Daring (Online). Metode: Self Directed Learning (SDL) dan Cooperative Learning (CoL). Penugasan: tugas kelompok mengulas topik dunia TI dalam IPTEK. & 1 x 3 x 50' tatap muka. & Mahasiswa mempelajari perkembangan IPTEK saat ini dan contoh perkembangan IPTEK dalam dunia pendidikan. & Kriteria: ketepatan dan penguasaan. Bentuk: keaktifan diskusi kelompok serta kesesuaian, kejelasan, dan kelugasan menjawab tugas. & Mahasiswa memiliki kemampuan memahami IPTEK serta dampak IPTEK dan solusinya. & 2 \\
4 & SCPMK0901-00201 & Pengertian etika, penggunaan etika dalam teknologi TI, jenis isu etika TI, peran etika bidang TI, serta etika dan tanggung jawab profesi TI. & Bentuk: Daring (Online). Metode: Self Directed Learning (SDL) dan Cooperative Learning (CoL). Penugasan: tugas kelompok mengulas standard internasional etika profesi IT. & 1 x 3 x 50' tatap muka. & Mahasiswa mempelajari contoh etika profesi dalam dunia IT dan memahami etika tak tertulis dalam dunia IT. & Kriteria: ketepatan dan penguasaan. Bentuk: keaktifan diskusi kelompok serta kesesuaian, kejelasan, dan kelugasan menjawab tugas. & Mahasiswa memiliki kemampuan memahami dan menjelaskan pengertian etika dalam penggunaan teknologi TI. & 2 \\
5 & SCPMK0901-00201 & Kuis 1: konsep teknologi, inovasi teknologi, perkembangan IPTEK, dan etika rekayasa. & Bentuk: Ujian Online. & 1 x 60' ujian. & Mahasiswa mengerjakan kuis secara mandiri. & Kuis. Mengacu rubrik RTM. & Mahasiswa mampu mengetahui dan memahami konsep teknologi, inovasi teknologi, IPTEK, dan etika TI. & 12 \\
6 & SCPMK0901-00201 & Pengertian ICT, perkembangan ICT, manfaat ICT dan penerapannya, serta perbedaan ICT dan TIK. & Bentuk: Daring (Online). Metode: Contextual Teaching and Learning (CTL) dan Self Directed Learning (SDL). Penugasan: tugas individu merangkum ulasan perkembangan ICT. & 1 x 3 x 50' tatap muka. & Mahasiswa mempelajari manfaat ICT dalam berbagai sektor dan perbedaan ICT dengan TIK. & Kriteria: ketepatan dan penguasaan. Bentuk: keaktifan diskusi kelompok serta kesesuaian, kejelasan, dan kelugasan menjawab tugas. & Mahasiswa dapat memahami perkembangan ICT, penerapannya, dan perbedaan ICT maupun TIK. & 2 \\
7 & SCPMK0901-00202 & Konsep sistem komputer: struktur komputer, perangkat I/O, interkoneksi antar komponen, register, memory, pemroses (CPU), unit kendali (CU), ALU, dan BUS. & Bentuk: Daring (Online). Metode: Contextual Teaching and Learning (CTL) dan Self Directed Learning (SDL). Penugasan: tugas individu merangkum konsep sistem komputer dan mengulas komponen komputer pada PC/laptop yang dimiliki. & 1 x 3 x 50' tatap muka. & Mahasiswa mempelajari konsep komputer dan komponen-komponen dalam suatu komputer. & Kriteria: ketepatan dan penguasaan. Bentuk: keaktifan diskusi kelompok serta kesesuaian, kejelasan, dan kelugasan menjawab tugas. & Kemampuan memahami dan menjelaskan struktur komputer dan komponennya. & 2 \\
8 & SCPMK0901-00201, SCPMK0901-00202 & UTS: konsep teknologi, inovasi teknologi, perkembangan IPTEK, etika rekayasa, perkembangan ICT, sistem komputer, dan konsep sistem komputer. & Bentuk: Ujian Online. & 1 x 60' ujian. & Mahasiswa mengerjakan UTS secara mandiri. & UTS. Mengacu rubrik RTM. & Kemampuan memahami cara menerapkan konsep-konsep TI, IPTEK, dan sistem komputer. & 15 \\
9 & SCPMK0901-00202 & Elemen sistem komputer, arsitektur sistem komputer, dan komponen sistem komputer. & Bentuk: Daring (Online). Metode: Self Directed Learning (SDL) dan Cooperative Learning (CoL). Penugasan: tugas kelompok mengulas arsitektur komputer pada komputer yang dimiliki. & 1 x 3 x 50' tatap muka. & Mahasiswa mempelajari arsitektur dalam sebuah komputer dan fungsi-fungsi dari komponen komputer. & Kriteria: ketepatan dan penguasaan. Bentuk: keaktifan diskusi kelompok serta kesesuaian, kejelasan, dan kelugasan menjawab tugas. & Kemampuan memahami dan menjelaskan dasar sistem komputer dan komponen sistem komputer. & 2 \\
10 & SCPMK0901-00203 & Representasi data: sistem bilangan, aritmatika, jenis tipe data, teori bilangan, konversi bilangan, dan penyajian data. & Bentuk: Daring (Online). Metode: Contextual Teaching and Learning (CTL) dan Self Directed Learning (SDL). Penugasan: tugas individu mengerjakan soal contoh representasi data. & 1 x 3 x 50' tatap muka. & Mahasiswa mempelajari konsep representasi data dan konversi suatu bilangan komputer. & Kriteria: ketepatan dan penguasaan. Bentuk: keaktifan diskusi kelompok serta kesesuaian, kejelasan, dan kelugasan menjawab tugas. & Kemampuan memahami cara menerapkan konsep representasi data meliputi sistem bilangan, jenis tipe data, dan teori bilangan. & 2 \\
11 & SCPMK0901-00203 & Aljabar Boolean: dasar operasi logika, gerbang logika, ekspresi Boolean, hukum aljabar Boolean, fungsi Boolean, dan aplikasinya. & Bentuk: Daring (Online). Metode: Contextual Teaching and Learning (CTL) dan Self Directed Learning (SDL). Penugasan: tugas individu mengerjakan soal contoh aljabar Boolean. & 1 x 3 x 50' tatap muka. & Mahasiswa mempelajari konsep aljabar Boolean dan operasi logika yang ada pada komputer. & Kriteria: ketepatan dan penguasaan. Bentuk: keaktifan diskusi kelompok serta kesesuaian, kejelasan, dan kelugasan menjawab tugas. & Kemampuan memahami cara menerapkan konsep aljabar Boolean, hukum, dan aplikasinya. & 2 \\
12 & SCPMK0901-00203 & Konsep flowchart, jenis flowchart, simbol flowchart, dan studi kasus penggunaan flowchart. & Bentuk: Daring (Online). Metode: Contextual Teaching and Learning (CTL) dan Self Directed Learning (SDL). Penugasan: tugas individu mengerjakan soal contoh flowchart. & 1 x 3 x 50' tatap muka. & Mahasiswa mempelajari konsep flowchart dan menerapkan flowchart dalam sebuah kasus. & Kriteria: ketepatan dan penguasaan. Bentuk: keaktifan diskusi kelompok serta kesesuaian, kejelasan, dan kelugasan menjawab tugas. & Kemampuan memahami cara menerapkan konsep dan studi kasus flowchart. & 2 \\
13 & SCPMK0901-00203 & Kuis 2: representasi data, aljabar Boolean, dan flowchart. & Bentuk: Ujian Online. & 1 x 60' ujian. & Mahasiswa mengerjakan kuis secara mandiri. & Kuis. Mengacu rubrik RTM. & Kemampuan memahami konsep representasi data, aljabar Boolean, dan flowchart. & 12 \\
14 & SCPMK0208-00204, SCPMK0502-00205 & Konsep jaringan komputer, konsep dan pengertian internet, jenis jaringan komputer, internet dan intranet, topologi jaringan, perangkat jaringan, dan pengenalan aspek keamanan dasar (firewall, autentikasi, ancaman umum) pada jaringan, sistem operasi, dan platform cloud. & Bentuk: Daring (Online). Metode: Self Directed Learning (SDL) dan Cooperative Learning (CoL). Penugasan: tugas kelompok mengulas topik perangkat jaringan komputer dan aspek keamanan dasarnya. & 1 x 3 x 50' tatap muka. & Mahasiswa mempelajari konsep dasar jaringan komputer serta jenis dan fungsi dari jaringan komputer, termasuk mengenali aspek keamanan dasar yang relevan. & Kriteria: ketepatan dan penguasaan. Bentuk: keaktifan diskusi kelompok serta kesesuaian, kejelasan, dan kelugasan menjawab tugas. & Kemampuan memahami cara menerapkan konsep jaringan komputer dan internet beserta topologi, jenis jaringan, perangkat jaringan, dan aspek keamanan dasar. & 2 \\
15 & SCPMK0208-00204 & Konsep aplikasi TI, jenis aplikasi TI, serta fungsi dan peranan TI dalam kehidupan sehari-hari dan perusahaan. & Bentuk: Daring (Online). Metode: Self Directed Learning (SDL) dan Cooperative Learning (CoL). Penugasan: tugas kelompok mengulas topik aplikasi TI. & 1 x 3 x 50' tatap muka. & Mahasiswa mempelajari konsep aplikasi serta jenis dan fungsi aplikasi TI dalam kehidupan sehari-hari. & Kriteria: ketepatan dan penguasaan. Bentuk: keaktifan diskusi kelompok serta kesesuaian, kejelasan, dan kelugasan menjawab tugas. & Kemampuan memahami cara menerapkan konsep aplikasi TI dalam kehidupan sehari-hari. & 2 \\
16 & SCPMK0208-00204 & Pengertian sertifikasi dan jenis sertifikasi bidang TI. & Bentuk: Daring (Online). Metode: Self Directed Learning (SDL) dan Cooperative Learning (CoL). Penugasan: tugas kelompok mengulas sertifikasi-sertifikasi pada bidang IT. & 1 x 3 x 50' tatap muka. & Mahasiswa mempelajari pengertian dan manfaat sertifikasi dalam bidang IT serta jenis sertifikasi dalam dunia IT. & Kriteria: ketepatan dan penguasaan. Bentuk: keaktifan diskusi kelompok serta kesesuaian, kejelasan, dan kelugasan menjawab tugas. & Kemampuan memahami jenis dan macam sertifikasi TI. & 2 \\
17 & SCPMK0901-00203, SCPMK0208-00204, SCPMK0502-00205 & UAS: representasi data, aljabar Boolean, flowchart, jaringan komputer dan internet, aplikasi TI, sertifikasi bidang TI, dan aspek keamanan dasar sistem operasi/jaringan/cloud. & Bentuk: Ujian Online. & 1 x 60' ujian. & Mahasiswa mengerjakan UAS secara mandiri. & UAS. Mengacu rubrik RTM. & Kemampuan menerapkan konsep komputasi, jaringan komputer dan internet, aplikasi TI, sertifikasi bidang TI, dan pengenalan keamanan dasar. & 35 \\
\end{longtblr}
